\section{Ordinary-full source conversions and constructor redexes}
\label{sec:p07-ch88-ordinary-checking}

These are further source-rule constructions, not an identification of the
ordinary full calculus with its normal-type variant. Only the corrected
ordinary full rules on Sheet~\ref{sheet:p07-ch88-full} are used. In
particular an object declaration requires an actual derivation of its
annotation at \(\ast\); target typing alone does not license it.
Named expressions below use the already proved raw indexed correspondence.
The original ordinary rules are \cite[pp.~99--102]{CoquandHuet1988}.

\begin{hblemma}{thm:p07-ch88-ordinary-star}{proved=annex:p7.ch88.ordinary-star}{Literal-star rigidity and its generation consequences}
For any scoped source expression \(E\), if \(TE=_\beta\ast\), then
\(E\) is literally \(\ast\). Thus source equality cannot convert
\(\ast\) to a different expression. No ordinary source lambda has
type \(\ast\). A typing \(\Gamma\vdash(MN):\ast\) has actual
generation premises \(\Gamma\vdash M:[x:A]\ast\) and
\(\Gamma\vdash N:A\). This last assertion does not assert a missing
source sort certificate for arbitrary inferred \(A\).
\end{hblemma}

\begin{hbtheorem}{thm:p07-ch88-ordinary-context}{proved=annex:p7.ch88.ordinary-context}{Context conversion with actual source annotation certificates}
Suppose \(A,A'\) each have their actual annotation certificate in
\(\Gamma\): object annotations are typed \(\ast\), and logical
annotations are valid logical contexts. Suppose also
\(\Gamma\vdash A\cong A'\). For every ordinary-full validity,
typing or equality judgment in \(\Gamma[x:A];\Theta\), one can
construct the same named judgment in \(\Gamma[x:A'];\Theta\).
The target suffix is valid; its dependent annotations are not discarded
or normalized. Indexed decoding gives the actual corresponding judgments.
No beta-equivalence-to-source-equality reflection is a premise.
\end{hbtheorem}

\begin{hbtheorem}{thm:p07-ch88-ordinary-equal-substitution}{proved=annex:p7.ch88.ordinary-equal-substitution}{Equality under two source-equal substitutions}
Suppose \(\Gamma\vdash N:A\), \(\Gamma\vdash N':A\), and
\(\Gamma\vdash N\cong N'\). If
\(\Gamma[x:A];\Theta\vdash E:P\), or
\(\Gamma[x:A];\Theta\vdash E\) is a validity judgment, then
\[
 \Gamma;([N/x]\Theta)\vdash [N/x]E\cong[N'/x]E.
\]
The right expression is compared in the first substituted environment,
as required by the full product/lambda congruence rules. This theorem
uses full typed reflexivity and simultaneous application congruence;
it is not advertised for arbitrary proof arguments in the restricted
variant.
\end{hbtheorem}

\begin{hbdefinition}{def:p07-ch88-ordinary-product-headed}{Product-headed equality proofs}
A product-headed source equality proof has products at both endpoints
of every outer symmetry/transitivity node. Its outer leaves are product
congruence or reflexivity of a typed object product or valid logical
product. The annotation/body equality subproofs may be arbitrary full
source proofs. Outer lambda, application and beta generators are excluded.
This is a property of an actual finite proof, not an assumption that
every object-product equality can already be transformed into that form.
\end{hbdefinition}

\begin{hbtheorem}{thm:p07-ch88-ordinary-product-injection}{proved=annex:p7.ch88.ordinary-product-injection}{Actual source product injection for the headed regime}
Suppose \(A\) has its actual source annotation certificate in
\(\Gamma\), and a product-headed proof gives
\(\Gamma\vdash[x:A]P\cong[x:B]Q\). Then one can construct
\[
 \Gamma\vdash A\cong B,\qquad
 \Gamma[x:A]\vdash P\cong Q.
\]
No certificate for the final \(B\) is assumed. If both endpoint
products are logical-context terms, every source equality proof between
them is automatically product-headed. Hence the conclusion is
unconditional for that kind-product case. It is not a claimed general
injection theorem for object products with arbitrary beta excursions.
\end{hbtheorem}

\begin{hbtheorem}{thm:p07-ch88-ordinary-root}{proved=annex:p7.ch88.ordinary-root}{Root subject reduction at every ordinary constructor type}
Let \(R=((\lambda x:A)M)N\) and \(C=[N/x]M\).
If \(\Gamma\vdash R:\Delta\), where \(\Delta\) is a logical
context, then actual ordinary-full derivations give
\[
 \Gamma\vdash C:\Delta,\qquad\Gamma\vdash R\cong C.
\]
The generation argument constructs the formerly missing source domain
equality and casts the argument to the literal lambda annotation.
For an arbitrary result type, the same conclusion holds when the actual
function-type conversion chain exposed by generation has product-headed
equality proofs. With literal argument typing \(\Gamma\vdash N:A\)
and body typing \(\Gamma[x:A]\vdash M:P\), both endpoints always
have the actual principal type \([N/x]P\); any explicitly supplied
outer conversion chain can be reapplied to the contractum.
\end{hbtheorem}

\begin{hbtheorem}{thm:p07-ch88-ordinary-application}{proved=annex:p7.ch88.ordinary-application}{Dependent application types under an actual equal argument}
Suppose \(\Gamma[x:A]\vdash B:\ast\), or that logical \(B\)
is valid there; suppose \(\Gamma\vdash F:[x:A]B\) and the same
three actual argument judgments as in
\cref{thm:p07-ch88-ordinary-equal-substitution}. Then
\[
 \Gamma\vdash[N/x]B\cong[N'/x]B,\qquad
 \Gamma\vdash FN':[N/x]B.
\]
The first is a constructed source equality, not merely a target beta
comparison. This supplies the codomain-transport step once the two
arguments' source typings at the same domain have been obtained.
\end{hbtheorem}

\begin{hbtheorem}{thm:p07-ch88-ordinary-example}{proved=annex:p7.ch88.ordinary-example}{Nonnormal declarations and an annotation contraction with the same old type}
In the ordinary valid environment
\[
 \Gamma=A:\ast,\ F:[u:A]\ast,\ a:A,\ f:[u:A]Fu,
\]
put \(n=((\lambda v:A)v)a\), \(B=Fn\), \(B'=Fa\), and
\(T=[h:B]B\). Actual full-source derivations give
\[
 n:A,\quad n\cong a,\quad B:\ast,\quad B':\ast,
 \quad B\cong B',\quad fn:B,\quad fa:B.
\]
The nonnormal declaration \(\Gamma[h:B]\) is valid. Changing it
to \(\Gamma[h:B']\) preserves the old judgment \(h:B\).
Furthermore \(t=(\lambda h:B)h\) contracts through its annotation
to \(t'=(\lambda h:B')h\), and both have the actual old type \(T\)
with \(\Gamma\vdash t\cong t'\).
The no-conversion normal-type variant cannot type the old nonnormal
annotation or retain the nonnormal old type \(T\).
\end{hbtheorem}

\begin{proofboundary}{def:p07-ch88-ordinary-boundary}{The remaining ordinary proof-valued conversion case}
General ordinary-full type correctness, sorted-beta equality reflection
and compatible subject reduction are not asserted complete here.
Constructor-valued root reduction now has its missing source product
injection and actual cast. For a proof-valued redex, the function type
is an object product. Its equality proof can leave the product head
through typed beta/application generators. The headed-path construction
does not yet eliminate those excursions. Neither one-way CC typing,
normal-type checking nor these successful nonnormal examples proves
that elimination or an arbitrary-annotation checker.
\end{proofboundary}
